\chapter{Introducción y Objetivos del Anteproyecto}
\label{sec:introduccion}

El presente anteproyecto documenta la propuesta para el Trabajo Integrador correspondiente al cuarto año de la
carrera de Ingeniería Electrónica en la Universidad Tecnológica Nacional, Facultad Regional Córdoba (UTN-FRC).
El objetivo central consiste en el diseño, desarrollo y validación experimental de una máquina de control numérico
computarizado (CNC) de tres ejes (X, Y, Z), orientada a aplicaciones de fresado ligero y prototipado rápido de
circuitos impresos (PCB).

Una particularidad tecnológica del proyecto reside en la reutilización sustentable de la estructura mecánica y
elementos de transmisión (guías lineales, husillos y chasis de soporte) provenientes de impresoras fuera de uso.
Esta aproximación fomenta la economía circular y plantea un desafío de ingeniería electromecánica: adaptar sistemas
de potencia y control a mecánicas heterogéneas no estandarizadas.

\section{Integración Académica Interdisciplinaria}
\label{sec:integracion}

El trabajo integrador articula los conocimientos teóricos y prácticos adquiridos en tres asignaturas clave del
nivel:

\begin{itemize}
  \item \textbf{Electrónica Aplicada II:} Implementación de etapas analógicas de acondicionamiento de señal, sensado
    de corriente en derivadores (\textit{shunts}) mediante amplificadores operacionales y detección comparada de
    sobrecorriente para actuar como final de carrera electrónico inteligente (detección de atascamiento o
    \textit{stall detection}).
  \item \textbf{Técnicas Digitales II:} Diseño e implementación de la arquitectura del sistema embebido bajo el
    sistema operativo en tiempo real FreeRTOS, desarrollo del módulo interpretador de comandos de código G (G-code),
    gestión de periféricos de comunicación (SPI, I2C, UART), temporizadores por hardware para generación de impulsos y
    diseño de la interfaz hombre-máquina (HMI) local.
  \item \textbf{Máquinas e Instalaciones Eléctricas:} Diseño y dimensionamiento de los controladores de potencia
    modulares (un canal independiente por eje) para motores paso a paso bipolares/unipolares (control de velocidad,
    dirección y evaluación de micropasos por troceado) y del controlador de potencia PWM para el motor DC de la fresa
    (\textit{spindle}).
\end{itemize}

\section{Campos de Aplicación}
\label{sec:aplicacion}

La plataforma desarrollada posee aplicaciones concretas en diversos ámbitos:
\begin{enumerate}
  \item Prototipado rápido de placas de circuito impreso mediante ruteado directo por aislamiento (PCB milling).
  \item Grabado y mecanizado liviano de materiales no ferrosos (madera, acrílico, MDF, placas FR4).
  \item Plataforma didáctica modular para la experimentación en control de movimiento, robótica y sistemas en tiempo
    real.
\end{enumerate}

\section{Arquitectura General del Sistema}
\label{sec:bloques_general}

El sistema global se estructura mediante una arquitectura modular distribuida en bloques funcionales, tal como se
esquematiza en la figura~\ref{fig.ej1:bloques_general}.

\begin{figure}[!ht]
  \centering
  \begin{tikzpicture}[node distance=1.8cm and 1.8cm, >=stealth, thick,
  block/.style={draw, fill=blue!10, rectangle, minimum height=1.2cm, minimum width=2.8cm, text centered, rounded corners, align=center},
  powerblock/.style={draw, fill=orange!15, rectangle, minimum height=1.2cm, minimum width=2.8cm, text centered, rounded corners, align=center},
  sensorblock/.style={draw, fill=purple!15, rectangle, minimum height=1.2cm, minimum width=2.8cm, text centered, rounded corners, align=center},
  mechblock/.style={draw, fill=yellow!20, rectangle, minimum height=1.2cm, minimum width=2.8cm, text centered, rounded corners, align=center}]

  % Nodos de la fila superior (Control y HMI)
  \node [block] (mcu) {Controlador Central \\ (FreeRTOS / TD2)};
  \node [block, left=of mcu] (hmi) {Interfaz HMI \\ (Pantalla + Encoder)};
  \node [block, right=of mcu] (storage) {Almacenamiento \\ SD (SPI/SDIO)};

  % Nodos de la fila intermedia y inferior
  \node [powerblock, below=of mcu] (driver) {Controlador Modular \\ Potencia (MyIE)};
  \node [sensorblock, left=of driver] (opamp) {Sensado Corriente \\ OpAmps (EA2)};
  \node [mechblock, below=of driver] (cnc) {Chasis Impresora \\ (Motores PaP + Spindle)};

  % Conexiones superiores
  \draw [<->] (hmi) -- node[above, font=\small] {I2C / GPIO} (mcu);
  \draw [<->] (storage) -- node[above, font=\small] {SPI} (mcu);

  % Conexiones de potencia y sensado
  \draw [->] (mcu) -- node[right, font=\small] {STEP/DIR/PWM} (driver);
  \draw [->] (driver) -- node[above, font=\small] {Tensión Shunt} (opamp);
  \draw [->] (driver) -- node[right, font=\small] {Potencia} (cnc);

  % Retornos y realimentación (Líneas ortogonales)
  \draw [->] (opamp) |- node[near end, left, font=\small] {EXTI / ADC} (mcu);
  \draw [->, dashed] (cnc) -| node[near start, below, font=\small] {Carga Mecánica} (opamp);
\end{tikzpicture}

  \caption{Diagrama en bloques general del sistema CNC modular y su articulación por asignatura.}
  \label{fig.ej1:bloques_general}
\end{figure}

El controlador central ejecuta FreeRTOS y se encarga de recibir las instrucciones, calcular trayectorias y emitir las
señales de control hacia los módulos de potencia. Cada canal de potencia acciona la máquina eléctrica correspondiente y
ofrece un punto de muestreo para que la etapa de acondicionamiento analógico supervise el consumo de corriente en todo
momento.
